\newpage\handout[true]
{\textsc{Math 2106-B: Foundations of Mathematical Proof}}{Fall 2025}
{Homework 5}
{\textit{Prof.} \href{https://ceheitsch.github.io/webpage/}{Dr. Christine Heitsch}}{Due Date: October 2nd}
{Math 2106-B: HW5}
{\large
  \noindent
  \textbf{Name:} \HandoutAuthor \smallskip \\

  \noindent
  \textbf{GTID:} 903743506 \smallskip \\

  \noindent
  %%% List anyone with whom you discussed the problems
  \textbf{Collaborators:} None. This material reflects independent preparation. \smallskip \\

  \noindent
  %%% List all resources used OTHER than the textbook, lecture notes, 
  %%% webwork problems, ICA questions, and previous homework assignments.
  \textbf{Outside resources:} This work was informed by MATH 2106 course materials 
  (ICAs: in-class activities and WebWork contents) from \HandoutInstitution\ (the course textbook is 
  Chartrand \parencite{chartrand_mathematical_2018}).   
  Outside references specifically include
  textbooks written by 
  Grimaldi \parencite{grimaldi_discrete_2004} (used in MATH 3012) and 
  Rosen \parencite{rosen_discrete_2019} (used in CS 2050). \smallskip \\

  \noindent
  \textbf{Notice}: The answers contain cross-references throughout this document as clickable hyperlinks in most PDF viewers. However, the Gradescope PDF view might not support clickable hyperlinks.
}\newpage
% \printbibliography[
  heading=bibintoc,   
  title={References}     % change “References” text if needed
]
% `heading=bibnumbered` if you prefer it numbered chapter (in \tableofcontents)
% `heading=bibintoc`    if you prefer it unnumbered but still in the ToC\newpage
% The following table presents the standard logical equivalences from these sources together with 
the author's own (or the student, myself, Jaehoon Song) clarifications and labels where applicable.


\begin{table}[h!]
  \centering
  \renewcommand{\arraystretch}{1.2}
  \setlength{\tabcolsep}{8pt}
  \begin{tabular}{|p{5.0cm}|p{5.0cm}|p{5.2cm}|}
    \hline
    \rowcolor[HTML]{E0E0E0}
    \multicolumn{2}{|l|}{\textbf{Logical Equivalences}} & \textbf{Name} \\
    \hline
    $\,p \AND q \equiv q \AND p$ & $\,p \OR q \equiv q \OR p$ & Commutative Law \\
    \hline
    $\, (p \AND q) \AND r \equiv p \AND (q \AND r)$ & $\, (p \OR q) \OR r \equiv p \OR (q \OR r)$ & Associative Law \\
    \hline
    $\,p \AND (q \OR r) \equiv (p \AND q) \OR (p \AND r)$ & $\,p \OR (q \AND r) \equiv (p \OR q) \AND (p \OR r)$ & Distributive Law \\
    \hline
    $\,\overline{p \AND q} \equiv \sNOT{p} \OR \sNOT{q}$ & $\,\overline{p \OR q} \equiv \sNOT{p} \AND \sNOT{q}$ & DeMorgan's Law \\
    \hline
    \multicolumn{2}{|l|}{$\,\overline{\overline{p}} \equiv p$} & Double Negation law \\
    \hline
    $\,p \AND \boldsymbol{T} \equiv p$ & $\,p \OR \boldsymbol{F} \equiv p$ & Identity Law \\
    \hline
    $\,p \OR \boldsymbol{T} \equiv \boldsymbol{T}$ & $\,p \AND \boldsymbol{F} \equiv \boldsymbol{F}$ & Domination Law \\
    \hline
    $\,p \AND p \equiv p$ & $\,p \OR p \equiv p$ & Idempotent Law \\
    \hline
    \multicolumn{2}{|l|}{$\,p \THEN q \equiv \sNOT{p} \OR q$} & Implication Equivalence \\
    \hline
    \multicolumn{2}{|l|}{$\,\overline{p \THEN q} \equiv p \AND \sNOT{q}$} & Negation of Implication \\
    \hline
    \multicolumn{2}{|l|}{$\,p \THEN q \equiv \sNOT{q} \THEN \sNOT{p}$} & Contrapositive Equivalence \\
    \hline
    \multicolumn{2}{|l|}{$\,p \OR \sNOT{p} \equiv \boldsymbol{T}$} & Law of Tautology \\
    \hline
    \multicolumn{2}{|l|}{$\,p \AND \sNOT{p} \equiv \boldsymbol{F}$} & Law of Contradiction \\
    \hline
  \end{tabular}
\end{table}\newpage
% 
The following definitions are presented from standard mathematical sources (\textbf{Sets}).
Clarifications and labels have been added where applicable.

\begin{define}{Basic Set Concepts}\phantomsection\label{def:basic-set-concepts} % \hyperref[def:basic-set-concepts]{YOUR_CONTEXT}
  \begin{enumerate}
    \item \textbf{Universal Set:} A universal set $\setUNIV$ is a
    set that contains all objects under consideration in a particular context or problem.
    \item \textbf{Empty Set:} The empty set, denoted $\setNULL$,
    is the unique set that contains no elements.
    \item \textbf{Subset:} Let $A$ and $B$ be sets. $A$ is a
    subset of $B$ ($A \subseteq B$) if (and only if)
    for all $x \in A$, $x \in B$.
    \item \textbf{Set Equality:} Let $A$ and $B$ be sets. The sets $A$ and $B$ are
    equal ($A = B$) if $A \subseteq B$ and $B \subseteq A$.
    \item \textbf{Power Set:} Let $A$ be a set. The power set of $A$, denoted $\sPOWER{A}$ or $2^A$,
    is the set of all subsets of $A$. That is, $\mathcal{P}(A) = \{S\in\setUNIV : S \subseteq A\}$.
  \end{enumerate}
\end{define}
The following definitions extend the understanding of set operations
and properties.


\begin{define}{Set Operations}\phantomsection\label{def:set-operations} % \hyperref[def:set-operations]{YOUR_CONTEXT}
  \begin{enumerate}
    \item \textbf{Complementary Set:} Let $A$ be a subset of the universal set $\setUNIV$.
    The complement of $A$, denoted $A^c$ or $\overline{A}$, is the set of all elements
    in $\setUNIV$ that are not in $A$. That is, $A^c = \{x \in \setUNIV : x \notin A\}$.
    \item \textbf{Union:} Let $A$ and $B$ be sets. The union of $A$ and $B$,
    denoted $A \cup B$, is the set of all elements that are in $A$, or in $B$, or in both.
    That is, $A \cup B = \{x : x \in A \lor x \in B\}$.
    \item \textbf{Intersection:} Let $A$ and $B$ be sets. The intersection of $A$ and $B$,
    denoted $A \cap B$, is the set of all elements that are in both $A$ and $B$.
    That is, $A \cap B = \{x : x \in A \land x \in B\}$.
    \item \textbf{Set Difference:} Let $A$, $B$ be subsets of a universal set $\setUNIV$.
    The set difference $A \sMINUS B$ is defined to be the set
    $\{ x \in \setUNIV \LDIV x \in A \text{ and } x \notin B\}$. \vspace*{10px}\\
    \textit{It is sometimes called the relative complement of $B$ in $A$,
    and denoted in the textbook as $A - B$.}
    \item \textbf{Cartesian Product:} Let $A$ and $B$ be sets. The Cartesian product of
    $A$ and $B$, denoted $A \sTIMES B$, is the set consisting of all ordered pairs:
    \[
    A \sTIMES B = \{(a, b) \LDIV a \in A, b \in B\}.
    \]
  \end{enumerate}
\end{define}


\newpage

The following claims are given in the reference as well. Let $A,B\in U$, a universal set.


\begin{claim}{Subset of Union: $A \subseteq \sOR{A}{B}$}\phantomsection\label{claim:subset-union} % \hyperref[claim:subset-union]{YOUR_CONTEXT}
  \begin{subprove}[bluecolor]
      Let $x\in U$, assume $x\in A$. Then the statement \inlinequote{$x\in A$ or $x\in B$} is true.
      Hence, by the definition, $\sOR{A}{B}=\{x\in U\LDIV x\in A$ or $x\in B\}$, $x\in \sOR{A}{B}$.
      Thus, $A \subseteq \sOR{A}{B}$.
  \end{subprove}
\end{claim}

\begin{claim}{Empty Subset: $\setNULL \subseteq A$}\phantomsection\label{claim:empty-subset} % \hyperref[claim:empty-subset]{YOUR_CONTEXT}
  \begin{subprove}[bluecolor]
      Let $x\in U$, assume $x\in \setNULL$. But there are no elements in $\setNULL$, so
      $x\in \setNULL$ is false. Since the premise is false, the implication if $x\in \setNULL$,
      then $x\in A$ is true. Thus, $\setNULL \subseteq A$ by definition of \hyperref[def:basic-set-concepts]{subsets}.
  \end{subprove}
\end{claim}

\begin{claim}{Subset as Equality: $\sOR{A}{B}=A$ if and only if $B\subseteq A$}\phantomsection\label{claim:subset-equality} % \hyperref[claim:subset-equality]{YOUR_CONTEXT}
  \begin{subprove}[bluecolor]
      Let $x\in U$, first, we assume $B\subseteq A$, we have already shown
      $A \subseteq \sOR{A}{B}$ by the claim, \hyperref[claim:subset-union]{Subset of Union}. 
      To show $\sOR{A}{B} \subseteq A$, assume $x\in \sOR{A}{B}$.
      By definition, $x\in A$ or $x\in B$. Considering $x\in B$ and assumption
      $B \subseteq A$, $x\in B$ implies $x\in A$. Thus, $\sOR{A}{B} \subseteq A$.
      Hence, if $B\subseteq A$, then $\sOR{A}{B} = A$. \vspace*{10px}\\
      Now, we assume $\sOR{A}{B}=A$. To show $B\subseteq A$, we now assume $x\in B$. Since $x\in B$, by definition
      of union, $x\in \sOR{A}{B} = A$. Therefore, $B\subseteq A$.
  \end{subprove}
\end{claim}


%%%%%%%%%%%%%%%%%%% 9/8/25

Let $\setUNIV$ be a set. 
\begin{claim}{ICA6}\phantomsection\label{claim:ica6} % \hyperref[claim:ica6]{YOUR_CONTEXT}
  For all $A,B\in\setUNIV$, $$A\sTIMES B =\setNULL\THEN \group{A=\setNULL\OR B=\setNULL}$$
  \begin{subprove}[bluecolor]
    Let $A,B\in\setUNIV$. We prove $A\sTIMES B=\setNULL$. Assume it is not the case that 
    $A=\setNULL$ or $B=\setNULL$. Thus, we have $A=\setNULL$ and $B=\setNULL$. 
    Since $A\ne\setNULL$ by assumption, there exists an element $a\in A$. Likewise, there 
    is some $b\in B$ since $B\ne\setNULL$. Thus, $\group{a,b}\in\sBUILD{(x,y)}{x\in A,y\in B}$ 
    by definition. But then, $A\sTIMES B\ne\setNULL$ by definition. Hence, the claim holds. 
  \end{subprove}
\end{claim}


\begin{define}{Relatively Prime Integers (Coprime)}\phantomsection\label{def:relatively-prime} % \hyperref[def:relatively-prime]{YOUR_CONTEXT}
  Let $x,y \in \mathbb{Z}$. The integers $x$ and $y$ are said to be
  \textbf{relatively prime} if and only if
  $\gcd(x,y) = 1$.
\end{define}
\newpage
% The following definitions and axioms are presented from standard mathematical sources (\textbf{Number Theory}) together with the author's own (or the student, myself, Jaehoon Song) clarifications and labels where applicable.

\begin{define}{Integers}[def:integers][red]
  $\setZ$: the set of all integers
\end{define}

\begin{define}{Rational Numbers}[def:rational-numbers][red]
  $\setQ$: the set of all rational numbers
\end{define}

\begin{define}{Real Numbers}[def:real-numbers][red]
  $\setR$: the set of all real numbers
\end{define}

\begin{define}{Natural Numbers}[def:natural-numbers][red]
  $\setN$ or $\setZ_{>0}$: the set of all natural numbers which are the set of positive integers (not including $0$)
\end{define}

\begin{define}{Positive Real Numbers}[def:positive-real-numbers][red]
  $\setR_{>0}$: the set of all positive real numbers
\end{define}

\begin{define}{Non-negative Real Numbers}[def:non-negative-real-numbers][red]
  $\setR_{\geq 0}$: the set of all non-negative real numbers
\end{define}


\newpage

The following definitions and claims extend the understanding of integers.

Here are the ICA2 Definitions and Axioms as follows.

\begin{define}{Even Integer}[def:even-integer][red]
  An integer $n$ is \underbar{even} if there exists 
  (existential quantification logic $\exists$) an integer $k$ such that
  \[
  n = 2k.
  \]
\end{define}

\begin{define}{Odd Integer}[def:odd-integer][red]
  An integer $n$ is \underbar{odd} if 
  there exists some integer $k$ such that
  \[
  n = 2k + 1.
  \]
\end{define}

\begin{define}{Closure under Addition}[def:closure-addition][red]
  The sum of two integers is again an 
  integer. (Closure under addition)
\end{define}

\begin{define}{Closure under Multiplication}[def:closure-multiplication][red]
  The product of two integers is again 
  an integer. (Closure under multiplication)
\end{define}

\begin{define}{Odd if and only if Not Even}[def:odd-iff-not-even][red]
  An integer is odd if and only if it 
  is not even.
\end{define}

Here are the ICA3 Definitions and Axioms as follows.

\begin{define}{Universal Closure under Addition}[def:universal-closure-addition][red]
  The sum of two integers (Universal quantification for all, for every, for each) is again an integer. (Closure under addition)
\end{define}

\begin{define}{Universal Closure under Multiplication}[def:universal-closure-multiplication][red]
  The product of two integers (Universal quantification for all, for every, for each) is again an integer. (Closure under multiplication)
\end{define}

\begin{define}{Even Integer Square}[def:even-integer-square][red]
  An integer $n$ is even if and only if $n^2$ is even.
\end{define}

\newpage

The following claims are given in the reference as well.

\begin{define}{Sum of Even Integers}[claim:sum-even-integers][red]
  The sum of two even integers is again an even integer.
  \begin{subprove}[red]
    Let $x$ be an integer and $y$ be an integer ($\forall x,y \in \setZ$). 
    Suppose $x$ is even and $y$ is even. Then, by \nameref{def:even-integer}, there exists an integer 
    $k$ and an integer $j$ such that $x=2k$ and $y=2j$. Then, 
    $x+y=\left(2k\right)+\left(2j\right)=2\left(k+j\right)$. By \nameref{def:closure-addition}, since $k$ and 
    $j$ are integers, then so is $\left(k+j\right)$. Hence, by definition, $x+y$ is an 
    even integer.
  \end{subprove}
\end{define}

The following definitions extend the understanding of integers including divisibility.

\begin{define}{Divisibility}[def:divisibility][red]
  An integer $n$ divides an integer $m$, denoted $n \LDIV m$, if there 
  exists an integer $k$ such that $m=n k$. Otherwise, $n \nmid m$.
\end{define}

The following definitions extend the classification of real numbers into rational and irrational numbers.

\begin{define}{Rational Number}[def:rational-number][red]
  A real number $x$ is \underbar{rational} if there exist integers $a$ and $b$ with $b \neq 0$ such that $bx = a$.
  Otherwise, $x$ is \underbar{irrational}.

  \textbf{Note}: Compare the rational number definition with the even/odd axiom \inlinequote{odd iff not even} from \nameref{def:odd-iff-not-even}.
\end{define}

The following claims are given in the reference as well.

\begin{define}{Irrationality of $\sqrt{2}$}[claim:sqrt2-irrational][red]
  $\sqrt{2}$ is irrational.
  \begin{subprove}[red]
    Suppose $\sqrt{2}$ is rational. Then, by \nameref{def:rational-number}, there exists $a,b\in \setZ$ 
    with $b\ne 0$ such that $b\cdot\sqrt{2}=a$. We may further assume that any factors 
    common to $a$ and $b$ have already been removed.\\

    Then, squaring both sides, this yields $2\cdot b^2 = a^2$.
    Since products of integers are again integers by \nameref{def:closure-multiplication}, we have that $a^2,b^2\in \setZ$.
    Thus, by definition, $a^2$ is even. Therefore, by \nameref{def:even-integer-square}, $a$ is even. 
    Thus, there exists $c\in\setZ$ such that $a=2c$. Then
    $$
    2b^2=a^2={2c}^2=4c^2
    $$
    Then, $b^2=2c^2$. Then, as before, $b^2$ is even, so $b$ is even.
    However, $a$ and $b$ both being even contradicts the fact that any factors 
    common to $a$ and $b$ have already been removed.
  \end{subprove}
\end{define}






%%%%%%%%%%%%%%%%%%% New style (refactor model)
\begin{define}{Relatively Prime Integers (Coprime)}[def:relatively-prime][red]
  Let $x,y \in \mathbb{Z}$. The integers $x$ and $y$ are said to be
  \textbf{relatively prime} if and only if
  $\gcd(x,y) = 1$.
\end{define}\newpage






%%%%%%%%%%%%%%%%%%%%%%%%%%%%%%%%%%%%%%%%%%%%%%%%%%%%%%%%%%%%%%%%%%%%%%%
% References
%%%%%%%%%%%%%%%%%%%%%%%%%%%%%%%%%%%%%%%%%%%%%%%%%%%%%%%%%%%%%%%%%%%%%%%
\textbf{Remark}: The following definitions (or results) are already proved in class or on previous ICA, 
Wbwk, Hmwk assignments, and contents of the textbook.

\begin{define}{One-to-One Correspondence (Bijective Function)}[def:bijective][red]
  Let $f : A \to B$ be a function. $f$ is \textbf{bijective} if and only if it is both injective and surjective.
\end{define}

\begin{define}{Theorem (Existence of Inverse)}[thm:inverse-function][red]
  The inverse relation $f^{-1}$ is a function if and only if $f$ is both injective and surjective.
  \begin{subprove}
    We consider the inverse relation $f^{-1} \subseteq B \times A$ for the function 
    \[
    f : A \to B \quad \text{on nonempty sets $A,B$.}
    \]
    Assume $f$ is a bijection, i.e., both one-to-one and onto. Let $b \in B$. Since $f$ is onto, by definition, there exists $a \in A$ with $f(a) = b$. So $(a,b) \in f$ and $(b,a) \in f^{-1}$. Hence for all $b \in B$ there exists $a \in A$ such that $(b,a) \in f^{-1}$. 

    Now let $a, a' \in A$. Suppose $(b,a) \in f^{-1}$ and $(b,a') \in f^{-1}$. 
    Then $(a,b), (a',b) \in f$, so $f(a) = b$ and $f(a') = b$. 
    But $f$ is one-to-one by assumption, and so $a = a'$. 
    Hence we have shown that each $b \in B$ has at least one image under 
    the inverse relation $f^{-1}$ and no more than one. 
    Thus the inverse relation is a function whenever $f$ is a bijection.

    \medskip
    Conversely, \dots
  \end{subprove}
  % \begin{subprove}
  %   ($\Rightarrow$) Assume $f^{-1}$ is a function. To show $f$ is injective, let $a_1, a_2 \in A$ such that $f(a_1) = f(a_2)$. 
  %   Then, $(a_1, f(a_1)) \in f$ and $(a_2, f(a_2)) \in f$. Since $f^{-1}$ is a function, $(f(a_1), a_1) \in f^{-1}$ and $(f(a_2), a_2) \in f^{-1}$. 
  %   By the definition of a function, $a_1 = a_2$. Thus, $f$ is injective.
    
  %   To show $f$ is surjective, let $b \in B$. Since $f^{-1}$ is a function, there exists an $a \in A$ such that $(b, a) \in f^{-1}$. 
  %   By the definition of inverse relation, $(a, b) \in f$, so $f(a) = b$. Thus, $f$ is surjective.
    
  %   ($\Leftarrow$) Assume $f$ is both injective and surjective. To show $f^{-1}$ is a function, let $(b, a_1), (b, a_2) \in f^{-1}$. 
  %   Then, $(a_1, b), (a_2, b) \in f$. Since $f$ is injective, it follows that $a_1 = a_2$. Thus, $f^{-1}$ satisfies the definition of a function.
  % \end{subprove}

\end{define}






\newpage













%%%%%%%%%%%%%%%%%%%%%%%%%%%%%%%%%%%%%%%%%%%%%%%%%%%%%%%%%%%%%%%%%%%%%%%
% Problems
%%%%%%%%%%%%%%%%%%%%%%%%%%%%%%%%%%%%%%%%%%%%%%%%%%%%%%%%%%%%%%%%%%%%%%%
\begin{enumerate} 
  \item Let $X, Y$, and $Z$ be nonempty sets with $A, B \subseteq X$ and $C, D \subseteq Y$. 
  Let $f: X \rightarrow Y$ and $g: Y \rightarrow Z$ be functions. 
  Prove the following statements or give the required counterexample.
  \begin{enumerate}
    \item Prove that if $f$ and $g$ are one-to-one, then so is $g \circ f$.
    \begin{define}{One-to-One (Injective Function)}[def:injective][red]
      Let $f : A \to B$ be a function. $f$ is \textbf{injective} (one-to-one) if and only if for all $a_1, a_2 \in A$, if $f(a_1) = f(a_2)$, then $a_1 = a_2$.
    \end{define}
    \begin{subprove}
      Suppose $f$ and $g$ are \nameref{def:injective} for the sake of direct proof.\proofstep
      Let $x_1, x_2 \in A$ be arbitrary such that $(g \circ f)(x_1) = (g \circ f)(x_2)$.
      By the definition of function composition, it follows that $$g(f(x_1)) = g(f(x_2)).$$
      By assumption, the function $g$ is one-to-one, it follows that $f(x_1) = f(x_2)$.
      Furthermore, $f$ is also one-to-one, it follows that $x_1 = x_2$.\proofstep
      Therefore, the function $g \circ f$ is one-to-one by definition.
    \end{subprove}
    \item Prove that if $g \circ f$ is one-to-one, then so is $f$.
    \begin{subprove}
      Assume $g \circ f$ is \nameref{def:injective}. Suppose $f$ is not one-to-one 
      for the sake of contradiction. \proofstep
      Let $x_1, x_2 \in A$ be arbitrary. Since $f$ is not one-to-one, 
      there exist $x_1, x_2 \in A$ such that $x_1 \neq x_2$ and $f(x_1) = f(x_2)$. 
      By the definition of function composition and the assumption, $$(g \circ f)(x_1) = g(f(x_1)) = g(f(x_2)) = (g \circ f)(x_2).$$
      Then, $x_1 = x_2$ since $g \circ f$ is one-to-one, which contradicts the assumption, $x_1 \neq x_2$. Therefore, $f$ must be one-to-one.
    \end{subprove}
    \newpage
    \item Prove that if $f$ is onto and $g \circ f$ is one-to-one, then $g$ is one-to-one.
    \begin{define}{Onto (Surjective Function)}[def:surjective][red]
      Let $f : A \to B$ be a function. $f$ is \textbf{surjective} (onto) if and only if for every $b \in B$, there exists $a \in A$ such that $f(a) = b$.
    \end{define}
    \begin{subprove}
      Let $y_1, y_2 \in Y$ be arbitrary elements. 
      Assume $f$ is \nameref{def:surjective} and $g \circ f$ is \nameref{def:injective} 
      for the sake of direct proof.\proofstep
      Suppose $y_1, y_2 \in Y$ be arbitrary such that $g(y_1) = g(y_2)$.
      Then, there exist $x_1, x_2 \in X$ such that 
      $f(x_1) = y_1$ and $f(x_2) = y_2$ for $y_1, y_2 \in Y$ since $f$ is onto by assumption. 
      Thus, by substitution, $g(f(x_1)) = g(f(x_2))$.
      Moreover, it implies $x_1 = x_2$ since $g \circ f$ is one-to-one. 
      \begin{define}{Function}[def:function][red]
        Let $A$ and $B$ be nonempty sets. A \textbf{function} $f$ from $A$ to $B$, denoted $f : A \to B$, 
        is a relation, $f\in A\sTIMES B$ such that 
        \begin{enumerate}
          \item For every $a \in A$, there exists an element $b \in B$ such that $f(a) = b$.
          \item For every $a \in A$ and for all $b, c \in B$, if $f(a) = b$ and $f(a) = c$, then $b = c$.
        \end{enumerate}
        where the set $A$ is called the \textbf{domain} of $f$, 
        and the set $B$ is called the \textbf{codomain} of $f$.
        If $f(a)=b$, then $b$ is called the \textbf{image} of $a$ under $f$.
        The function maps $a$ to $b$, and $f$ is often called a \textbf{mapping}.
      \end{define}
      \begin{define}{Domain and Codomain}[def:domain-codomain][red]
        If $f : A \to B$ is a function, then $A$ is called the \textbf{domain} of $f$, and $B$ is called the \textbf{codomain} of $f$.
      \end{define}
      By the definition of \nameref{def:function}, there must be a mapping of $x$ to $y$ in \nameref{def:domain-codomain} such that
      $$y_1 = f(x_1) = f(x_2) = y_2.$$
      Thus, $g$ is one-to-one by definition.\proofstep
      Therefore, by direct proof, $f$ is onto and $g \circ f$ is one-to-one only if $g$ is one-to-one.
    \end{subprove}
    \newpage
    \item Give an example where $g \circ f$ is one-to-one, but $g$ is not.
    \begin{subanswer}
      Let $X=\{1,2\}, Y=\{3,4,5\}, Z=\{6,7\}$. 
      Let $f: X \rightarrow Y$ and $g: Y \rightarrow Z$ be functions defined as follows for the sake of counterexample.
      \[
        f(1)=3, \quad f(2)=4; \qquad g(3)=6, \quad g(4)=7, \quad g(5)=7.
      \]
      % Diagram of f and g and g o f
      diagramed as follows. \label{prob:1d}
      \[
        \xymatrix@R=2em@C=2em{
          X \xyArc[r]^{f}   & Y \xyArc[r]^{g}   & Z \\
          \{1\} \xyArc[r]   & \{3\} \xyArc[r]   & \{6\} \\
          \{2\} \xyArc[r]   & \{4\} \xyArc[r]   & \{7\} \\
          {}                & \{5\} \xyArc[ur]  & {} \\
        }
      \]
      Then, the composition $g \circ f$ is one-to-one since 
      $(g \circ f)(1) = g(f(1)) = g(3) = 6$ and $(g \circ f)(2) = g(f(2)) = g(4) = 7$ which follows the definition of \nameref{def:injective}.
      However, $g$ is not one-to-one since $g(4) = 7$ and $g(5) = 7$, but $4 \neq 5$.\proofstep
      Therefore, by counterexample, there exists a case that $g \circ f$ is one-to-one, but $g$ is not.
    \end{subanswer}
  \end{enumerate}



  \newpage
  \item Let $X, Y$, and $Z$ be nonempty sets with $A, B \subseteq X$ and $C, D \subseteq Y$. 
  Let $f: X \rightarrow Y$ and $g: Y \rightarrow Z$ be functions. 
  Prove the following statements or give the required counterexample.
  \begin{enumerate}
    \item Prove that if $f$ and $g$ are onto, then so is $g \circ f$.
    \begin{subprove}
      Let $z \in Z$ be an arbitrary image for $g \circ f$. 
      Then, assume $f$ and $g$ are onto for the sake of direct proof.
      By definition of \nameref{def:surjective}, 
      there exists $x \in X$ such that $f(x) = y$ since $f$ is onto.
      Also, there exists $y \in Y$ such that $g(y) = z$ since $g$ is onto.\proofstep
      Thus, for any $z \in Z$, there exists $x \in X$ such that $(g \circ f)(x) = z$.
      Therefore, $g \circ f$ is onto by definition.
    \end{subprove}
    \item Prove that if $g \circ f$ is onto, then so is $g$.
    \begin{subprove}
      Let $z \in Z$ be an arbitrary image for $g \circ f$. 
      Assume $g \circ f$ is onto for the sake of direct proof.
      By definition of \nameref{def:surjective}, 
      there exists $x \in X$ such that $(g \circ f)(x) = z$.
      Moreover, there exists $y \in Y$ such that $f(x) = y$ since $f$ is a \nameref{def:function}.
      By substitution, $g(y) = z$.\proofstep
      Thus, for any $z \in Z$, there exists $y \in Y$ such that $g(y) = z$.
      Therefore, $g$ is onto by definition.
    \end{subprove}
    \item Prove that if $g$ is one-to-one and $g \circ f$ is onto, then $f$ is onto.
    \begin{subprove}
      Let $y \in Y$ be an arbitrary image for $f$. 
      Assume $g$ is \nameref{def:injective} and $g \circ f$ is \nameref{def:surjective} for the sake of direct proof.
      By definition, there exists $z \in Z$ such that $g(y) = z$ since $g$ is a \nameref{def:function}.
      Moreover, there exists $x \in X$ such that $(g \circ f)(x) = z$ for the $z$ since $g \circ f$ is onto.
      It follows that 
      $$g(f(x)) = g(y) = z$$
      which implies $f(x) = y$ since $g$ is \nameref{def:injective}.\proofstep
      Thus, for any $y \in Y$, there exists $x \in X$ such that $f(x) = y$.
      Therefore, $f$ is onto by definition.
    \end{subprove}
    \item Give an example where $g \circ f$ is onto, but $f$ is not.
    \begin{subanswer}
      Let $X=\{1,2\}, Y=\{3,4,5\}, Z=\{6,7\}$. 
      Let $f: X \rightarrow Y$ and $g: Y \rightarrow Z$ be functions defined as follows for the sake of counterexample.
      \[
        f(1)=3, \quad f(2)=4; \qquad g(3)=6, \quad g(4)=7, \quad g(5)=7.
      \]
      It has been already diagramed in the previous problem, \parenref{prob:1d}.
      Then, the composition $g \circ f$ is onto since 
      $(g \circ f)(1) = g(f(1)) = g(3) = 6$ and $(g \circ f)(2) = g(f(2)) = g(4) = 7$ which follows the definition of \nameref{def:surjective}.
      However, $f$ is not onto since there is no $x \in X$ such that $f(x) = 5$.\proofstep
      Therefore, by counterexample, there exists a case that $g \circ f$ is onto, but $f$ is not.
    \end{subanswer}
  \end{enumerate}


  \newpage 
  \item Let $X, Y$, and $Z$ be nonempty sets with $A, B \subseteq X$ and $C, D \subseteq Y$. 
  Let $f: X \rightarrow Y$ and $g: Y \rightarrow Z$ be functions. 
  Prove the following statements.
  \begin{define}{Set Equality}[def:set-equality][red]
    Let $A$ and $B$ be sets. The sets $A$ and $B$ are equal ($A = B$) if $A \subseteq B$ and $B \subseteq A$.
  \end{define}
  \begin{define}{Subset}[def:subset][red]
    Let $A$ and $B$ be sets. $A$ is a subset of $B$ ($A \subseteq B$) if (and only if) for all $x \in A$, $x \in B$.
  \end{define}
  \textbf{Remark}: Let $A$ and $B$ be arbitrary sets. By definition of \nameref{def:set-equality}, 
  $A = B$ requires showing both directions of \nameref{def:subset}s, $A \subseteq B$ and $B \subseteq A$.
  \begin{define}{Image and Inverse Image}[def:image][red]
    Let $f : X \to Y$ be a function, $C\subseteq X$, and $D\subseteq Y$. 
    \begin{enumerate}
      \item The \textbf{image} $f(C)$ is the set $\{f(c) \in Y : c \in C\} \subseteq Y$.
      \item The \textbf{inverse image} $f^{-1}(D)$ is the set $\{x \in X : f(x) \in D\} \subseteq X$.
    \end{enumerate}
  \end{define} 
  \begin{enumerate}
    \item $f(A \sOR B)=f(A) \sOR f(B)$ \label{prob:3a}
    \begin{subprove}
      ($\subseteq$) Let $y \in f(A \sOR B)$ be arbitrary. Then, by definition of \nameref{def:image}, 
      there exists $x \in A \sOR B$ such that $f(x) = y$.
      \begin{define}{Union}[def:union][red]
        Let $A$ and $B$ be sets. The union of $A$ and $B$, denoted $A \sOR B$, is the set of all elements that are in $A$, or in $B$, or in both. That is, $A \sOR B = \{x : x \in A \lor x \in B\}$.
      \end{define}
      Since $x \in A \sOR B$, by definition of \nameref{def:union}, $x \in A$ or $x \in B$. By the cases as follows,
      \begin{itemize}
        \item If $x \in A$, then $y = f(x) \in f(A)$ by definition.
        \item If $x \in B$, then $y = f(x) \in f(B)$ by definition.
      \end{itemize}
      which logically equivalent to $y \in f(A) \sOR f(B)$.
      Thus, $f(A \sOR B) \subseteq f(A) \sOR f(B)$.\proofstep
      ($\supseteq$) Let $y \in f(A) \sOR f(B)$ be arbitrary. Then, by definition (trivially), 
      $y \in f(A)$ or $y \in f(B)$.
      \begin{define}{Subset of Union}[claim:subset-union][red]
        Let $A$ and $B$ be sets. Then, $A \subseteq \sOR{A}{B}$
        \begin{subprove}[red]
          Let $x\in U$, assume $x\in A$. Then the statement \inlinequote{$x\in A$ or $x\in B$} is true.
          Hence, by definition of \nameref{def:union}, $\sOR{A}{B}=\{x\in U\LDIV x\in A$ or $x\in B\}$, $x\in \sOR{A}{B}$.
          Thus, $A \subseteq \sOR{A}{B}$.
        \end{subprove}
      \end{define}
      \begin{itemize}
        \item If $y \in f(A)$, then there exists $x \in A$ such that $f(x) = y$ by definition.
        Since $x \in A$, it follows that $x \in A \sOR B$ by \nameref{claim:subset-union}.
        \item If $y \in f(B)$, then there exists $x \in B$ such that $f(x) = y$ by definition.
        Since $x \in B$, it follows that $x \in A \sOR B$ by \nameref{claim:subset-union}.
      \end{itemize}
      In both cases, $y \in f(A \sOR B)$.
      Thus, $f(A) \sOR f(B) \subseteq f(A \sOR B)$.\proofstep
      Therefore, by ($\subseteq$) and ($\supseteq$), $f(A \sOR B)=f(A) \sOR f(B)$.
    \end{subprove}

    \newpage
    \item $f^{-1}(C \sOR D)=f^{-1}(C) \sOR f^{-1}(D)$
    \begin{subprove}
      ($\subseteq$) Let $x \in f^{-1}(C \sOR D)$ be arbitrary. Then, by definition of \nameref{def:image}, 
      $f(x) \in C \sOR D$.
      Since $f(x) \in C \sOR D$, by definition of \nameref{def:union}, $f(x) \in C$ or $f(x) \in D$.
      \begin{itemize}
        \item If $f(x) \in C$, then $x \in f^{-1}(C)$ by definition.
        \item If $f(x) \in D$, then $x \in f^{-1}(D)$ by definition.
      \end{itemize}
      which logically equivalent to $x \in f^{-1}(C) \sOR f^{-1}(D)$.
      Thus, $f^{-1}(C \sOR D) \subseteq f^{-1}(C) \sOR f^{-1}(D)$.\proofstep
      ($\supseteq$) Let $x \in f^{-1}(C) \sOR f^{-1}(D)$ be arbitrary. Then, by definition (trivially), 
      $x \in f^{-1}(C)$ or $x \in f^{-1}(D)$.
      \begin{itemize}
        \item If $x \in f^{-1}(C)$, then $f(x) \in C$ by definition.
        Since $f(x) \in C$, it follows that $f(x) \in C \sOR D$ by \nameref{claim:subset-union}.
        \item If $x \in f^{-1}(D)$, then $f(x) \in D$ by definition.
        Since $f(x) \in D$, it follows that $f(x) \in C \sOR D$ by \nameref{claim:subset-union}.
      \end{itemize}
      In both cases, $x \in f^{-1}(C \sOR D)$.
      Thus, $f^{-1}(C) \sOR f^{-1}(D) \subseteq f^{-1}(C \sOR D)$.\proofstep
      Therefore, by ($\subseteq$) and ($\supseteq$), $f^{-1}(C \sOR D)=f^{-1}(C) \sOR f^{-1}(D)$.
    \end{subprove}
    \item $f^{-1}(C \sAND D)=f^{-1}(C) \sAND f^{-1}(D)$
    \begin{subprove}
      ($\subseteq$) Let $x \in f^{-1}(C \sAND D)$ be arbitrary. Then, by definition of \nameref{def:image}, 
      $f(x) \in C \sAND D$.
      \begin{define}{Intersection}[def:intersection][red]
        Let $A$ and $B$ be sets. The intersection of $A$ and $B$, denoted $A \sAND B$, is the set of all elements that are in both $A$ and $B$. That is, $A \sAND B = \{x : x \in A \land x \in B\}$.
      \end{define}
      Since $f(x) \in C \sAND D$, by definition of \nameref{def:intersection}, $f(x) \in C$ and $f(x) \in D$.
      Therefore, $x \in f^{-1}(C)$ and $x \in f^{-1}(D)$ by definition,
      which logically equivalent to $x \in f^{-1}(C) \sAND f^{-1}(D)$.
      Thus, $f^{-1}(C \sAND D) \subseteq f^{-1}(C) \sAND f^{-1}(D)$.\proofstep
      ($\supseteq$) Let $x \in f^{-1}(C) \sAND f^{-1}(D)$ be arbitrary. Then, by definition of \nameref{def:intersection}, 
      $x \in f^{-1}(C)$ and $x \in f^{-1}(D)$.
      Since $x \in f^{-1}(C)$, it follows that $f(x) \in C$
      and $x \in f^{-1}(D)$, it follows that $f(x) \in D$ by definition,
      which logically equivalent to $f(x) \in C \sAND D$.
      Furthermore, $x \in f^{-1}(C \sAND D)$.
      Thus, $f^{-1}(C) \sAND f^{-1}(D) \subseteq f^{-1}(C \sAND D)$.\proofstep
      Therefore, by ($\subseteq$) and ($\supseteq$), $f^{-1}(C \sAND D)=f^{-1}(C) \sAND f^{-1}(D)$.
    \end{subprove}
  \end{enumerate}







  \newpage
  \item Let $X, Y$, and $Z$ be nonempty sets with $A, B \subseteq X$ and $C, D \subseteq Y$. 
  Let $f: X \rightarrow Y$ and $g: Y \rightarrow Z$ be functions. 
  Prove the following statements.
  \begin{enumerate}
    \item $f(A \sAND B) \subseteq f(A) \sAND f(B)$. \label{prob:4a}
    \begin{subprove}
      Let $y \in f(A \sAND B)$ be arbitrary for the sake of directly showing ($\subseteq$). 
      Then, by definition of \nameref{def:image}, there exists $x \in A \sAND B$ such that $f(x) = y$.
      Since $x \in A \sAND B$, by definition of \nameref{def:intersection}, $x \in A$ and $x \in B$.
      Since $x \in A$ and $f(x) = y$, it follows that $y \in f(A)$ by definition.
      Since $x \in B$ and $f(x) = y$, it follows that $y \in f(B)$ by definition.
      Thus, $y \in f(A) \sAND f(B)$ by definition of \nameref{def:intersection}.
      Therefore, $f(A \sAND B) \subseteq f(A) \sAND f(B)$.
    \end{subprove}
    \item $f$ is one-to-one if and only if $f(A \sAND B)=f(A) \sAND f(B)$ for all $A, B \subseteq X$.
    \begin{subprove}
      ($\Rightarrow$) Suppose $f$ is one-to-one for the sake of direct proof. 
      To show $f(A \sAND B)=f(A) \sAND f(B)$ for all $A, B \subseteq X$, 
      both ($\subseteq$) and ($\supseteq$) are required.
      Fortunately, ($\subseteq$) direction was already proved in part \parenref{prob:4a}.\proofstep
      For ($\supseteq$), let $y \in f(A) \sAND f(B)$ be arbitrary. 
      Then, by definition of \nameref{def:intersection}, $y \in f(A)$ and $y \in f(B)$.
      Since $y \in f(A)$, there exists $x_1 \in A$ such that $f(x_1) = y$ by definition of \nameref{def:image}.
      Since $y \in f(B)$, there exists $x_2 \in B$ such that $f(x_2) = y$ by definition.
      Since $f$ is \nameref{def:injective} by assumption and $f(x_1) = f(x_2) = y$, it follows that $x=x_1 = x_2$.
      Thus, $x \in A \sAND B$ which is equivalent to $y = f(x_1) \in f(A \sAND B)$.
      Therefore, $f(A) \sAND f(B) \subseteq f(A \sAND B)$.\proofstep
      Thus, $f$ is one-to-one only if $f(A \sAND B)=f(A) \sAND f(B)$ for all $A, B \subseteq X$.\proofstep
      ($\Leftarrow$) Let $A = \{x_1\} \subseteq X$ and $B = \{x_2\} \subseteq X$ such that $f(x_1) = f(x_2)$.
      Suppose $f(A \sAND B)=f(A) \sAND f(B)$ for all $A, B \subseteq X$ for the sake of direct proof.
      By definition of \nameref{def:domain-codomain},
      \begin{enumerate}
        \item $f(A) = \{f(x_1)\}$
        \item $f(B) = \{f(x_2)\} = \{f(x_1)\}$
      \end{enumerate}
      Since $f(x_1) = f(x_2)$. Moreover,
      $$f(A \sAND B) = f(A) \sAND f(B) = \{f(x_1)\}$$ by assumption. 
      Also, by definition of \nameref{def:domain-codomain}, $x_1 \in A \sAND B$.
      Here, shortly claim the following.
      \begin{define}{Subset of Intersection}[claim:subset-intersection][red]
        Let $A$ and $B$ be sets. Then, $A \sAND B \subseteq B$ and $A \sAND B \subseteq A$
        \begin{subprove}[red]
          Let $x\in \setUNIV$, assume $x\in A \sAND B$. Then the statement \inlinequote{$x\in A$ and $x\in B$} is true.
          Hence, by definition of \nameref{def:intersection}, $\sAND{A}{B}=\{x\in U\LDIV x\in A$ and $x\in B\}$, $x\in A$ and $x\in B$.
          Thus, $A \sAND B \subseteq A$ and $A \sAND B \subseteq B$.
        \end{subprove}
      \end{define}
      By \nameref{claim:subset-intersection}, $x_1 \in A=\set{x_1}$ (trivial) 
      and $x_1 \in B=\set{x_2}$, which implies $x_1 = x_2$.\proofstep
      Thus, $f$ is one-to-one if $f(A \sAND B)=f(A) \sAND f(B)$ for all $A, B \subseteq X$.\proofstep
      Therefore, by ($\Rightarrow$) and ($\Leftarrow$), $f$ is one-to-one 
      if and only if $f(A \sAND B)=f(A) \sAND f(B)$ for all $A, B \subseteq X$.
    \end{subprove}
  \end{enumerate}
\end{enumerate}
